\documentclass[conference]{IEEEtran}

\usepackage[utf8]{inputenc}
\usepackage[T1]{fontenc}
\usepackage{amsmath,amssymb,amsthm}
\usepackage{graphicx}
\usepackage{booktabs}
\usepackage{multirow}
\usepackage{url}
\usepackage{hyperref}
\usepackage{xcolor}
\usepackage{caption}
\usepackage{subcaption}
\usepackage{array}
\usepackage{algorithm}
\usepackage{algpseudocode}

\newcolumntype{C}[1]{>{\centering\arraybackslash}p{#1}}
\newtheorem{definition}{Definition}
\newtheorem{proposition}{Proposition}
\newtheorem{theorem}{Theorem}

\begin{document}

\title{The Information-Theoretic Cost of Hierarchy:\\
A Channel-Theoretic Analysis of LLM Multi-Agent Organizations}

\author{
\IEEEauthorblockN{Koh Shinoda}
\IEEEauthorblockA{
% Affiliation here
}
}

\maketitle

% ============================================================
\begin{abstract}
% ============================================================

% [STATUS: ドラフト — 実験完了後に数値を埋める]

LLM multi-agent systems adopt diverse organizational structures---hierarchies, review protocols, communication topologies---yet no theoretical framework explains \emph{why} some structures outperform others.
Prior empirical work (including our OrgBench benchmark) has shown that shallow hierarchies consistently outperform deep ones, but the causal mechanism remains unknown.
%
We address this gap by modeling multi-agent organizations as \textbf{information processing graphs}, where each agent acts as a communication channel with finite capacity.
Applying the Data Processing Inequality (DPI), we derive theoretical predictions:
(1)~information fidelity decays monotonically with chain depth,
(2)~mesh topologies preserve information better than chains through path redundancy, and
(3)~review gates function as feedback channels that reduce error rates without increasing capacity.
%
However, our empirical analysis reveals that LLM agents are \emph{not} simple noisy channels:
they exhibit three distinct modes---\textbf{preservation} (relay), \textbf{compression} (filtering), and \textbf{augmentation} (knowledge injection)---depending on their role and input volume.
We identify \textbf{channel saturation} as a key phenomenon: when input exceeds an agent's effective capacity, output quality collapses (compression ratio drops to 0.06).
%
Through controlled experiments with [TODO: N] runs across [TODO: M] configurations, we validate these predictions using six information-theoretic metrics including tracer survival analysis and semantic similarity decay.
Our framework provides the first principled design guidelines for multi-agent topology: optimize for the \textbf{minimum-capacity bottleneck} in the information flow graph, not the total path length.

% [TODO: 実験結果確定後に具体的数値を追加]
\end{abstract}

\begin{IEEEkeywords}
Information theory, Multi-agent systems, Channel capacity, Data processing inequality, Organizational design, Large language models
\end{IEEEkeywords}


% ============================================================
\section{Introduction}
\label{sec:intro}
% ============================================================

% --- 1段落目: 実務的問題 ---

The rapid proliferation of LLM-based multi-agent frameworks---AutoGen~\cite{wu2023autogen}, MetaGPT~\cite{hong2023metagpt}, ChatDev~\cite{qian2024chatdev}, CrewAI, LangGraph---has made organizational design a daily engineering decision.
Practitioners must choose hierarchy depth, communication topology, and review protocols, yet no principled theory guides these choices.
The current state of practice is trial-and-error: developers test configurations empirically, at costs of \$100--1,000 per project, without understanding \emph{why} one structure outperforms another.

% --- 2段落目: 学術的ギャップ ---

Recent work has begun to treat topology as an experimental variable.
MacNet~\cite{qian2025macnet} demonstrated that irregular DAG topologies outperform regular ones at ICLR 2025;
G-Designer~\cite{yao2025gdesigner} used graph neural networks to generate task-adaptive topologies at ICML 2025.
However, all such studies share a fundamental limitation:
they show \emph{what} works but not \emph{why}.
MacNet cannot explain why irregular topologies are superior.
G-Designer generates effective topologies but cannot predict their performance without running the full pipeline.
Without a theoretical framework, each empirical finding remains an isolated data point that cannot generalize to new conditions (different models, agent counts, or task types).

% --- 3段落目: OrgBenchの知見と「なぜ」 ---

Our prior work, OrgBench~\cite{shinoda2026orgbench}, systematically compared 11 information flow designs through 275 controlled runs.
The results were clear: shallow hierarchies outperform deep ones ($\Delta = 0.38$--$0.46$, $\eta^2 = 0.21$), mesh communication provides marginal gains over hub-and-spoke, and review protocols improve quality modestly ($\Delta = 0.24$).
But \emph{why} does depth hurt?
At least three hypotheses are plausible:
(a)~information distortion accumulates across relay steps;
(b)~intermediate layers over-filter, discarding relevant content;
(c)~intermediate layers consume token budget that could improve the final output.
Distinguishing these requires a theoretical framework---one that can model information flow, capacity constraints, and distortion accumulation across agent chains.

% --- 4段落目: 情報理論の提案 ---

We propose that \textbf{information theory} provides the natural framework.
Multi-agent organizations are, literally, information processing systems: agents receive text, process it, and produce text.
Each agent can be modeled as a communication channel with measurable input, output, and capacity constraints.
Shannon's foundational concepts---mutual information, channel capacity, the Data Processing Inequality, and rate-distortion theory---apply directly.

% [TODO: 予備分析の数値を具体的に引用]

Indeed, preliminary analysis of our OrgBench data reveals striking information-theoretic phenomena.
The Team Lead agent consistently outputs only 35\% of its input tokens (compression ratio 0.35), acting as a \emph{fixed-rate information filter}.
In mesh configurations, the CFO and CDO agents receive 4$\times$ more input than in hub configurations, and their compression ratio collapses to 0.06---they output only 6\% of input, exhibiting \textbf{channel saturation}.
Meanwhile, the CEO agent consistently produces \emph{more} tokens than it receives (compression ratio 1.85), functioning as an \textbf{information amplifier} that augments input with learned knowledge.

These observations suggest that LLM agents are not simple noisy channels subject to the classical Data Processing Inequality, but \emph{adaptive channels} whose behavior varies with input volume and role---a phenomenon that demands a new theoretical treatment.

% --- 5段落目: 貢献 ---

Our contributions are:
\begin{enumerate}
    \item A \textbf{formal model} of multi-agent organizations as information processing graphs, where each agent is a communication channel with role-dependent capacity and three operating modes (preservation, compression, augmentation).
    \item \textbf{Theoretical predictions} derived from the Data Processing Inequality, including the Depth-Distortion Principle (information fidelity decays exponentially with critical-path depth) and the decomposition of information loss into intentional filtering ($D_k^{\text{filter}}$) and unintentional noise ($D_k^{\text{noise}}$).
    \item \textbf{Six information-theoretic metrics} for quantifying information flow in multi-agent systems: compression ratio, semantic similarity decay, entity preservation rate, numerical fidelity, information density, and tracer survival rate.
    \item \textbf{Empirical validation} through [TODO: N] controlled experiments, demonstrating that information-theoretic metrics predict output quality better than structural metrics (graph degree, path length) and identifying channel saturation as the primary mechanism behind the ``depth penalty.''
\end{enumerate}


% ============================================================
\section{Related Work}
\label{sec:related}
% ============================================================

% --- 2.1: マルチエージェントのトポロジー研究 ---
\subsection{Topology as a Variable in Multi-Agent Systems}

The treatment of communication topology as an experimental variable has evolved rapidly.
First-generation frameworks (AutoGen~\cite{wu2023autogen}, MetaGPT~\cite{hong2023metagpt}, ChatDev~\cite{qian2024chatdev}) embedded fixed topologies without comparison.
Second-generation work introduced topology variation:
MacNet~\cite{qian2025macnet} (ICLR 2025) explored DAG-based topologies and discovered a logistic scaling law relating agent count to collective performance;
G-Designer~\cite{yao2025gdesigner} (ICML 2025) used variational graph autoencoders for task-adaptive topology generation;
MultiAgentBench~\cite{multiagentbench2025} (ACL 2025) compared star, chain, tree, and graph topologies.

Concurrently, AgentPrune~\cite{agentprune2025} (ICLR 2025) learns sparse communication graphs, showing that pruned topologies outperform dense ones while reducing cost---an empirical finding our saturation theory explains (dense routing saturates weak agents).
AgentDropout~\cite{agentdropout2025} (ACL 2025) dynamically eliminates redundant agents and edges.
Shen et al.~\cite{shen2025infoprop} directly study information propagation effects of topologies in LLM-MAS, the closest prior work to ours.
Cemri et al.~\cite{cemri2025fail} (NeurIPS 2025 spotlight) analyze 1,600+ failure traces across 7 MAS frameworks and find that 79\% of failures stem from coordination issues---supporting the thesis that information flow design, not compute, is the primary bottleneck.

All these studies characterize topology effects empirically.
Our work differs by providing a \textbf{structured information-flow analysis} that decomposes the effects into distortion ($D_k$) and augmentation ($A_k$) at each agent step, enabling predictions for unseen configurations.

% --- 2.2: 情報理論とLLM ---
\subsection{Information Theory and Language Models}

The connection between language modeling and information theory has deep roots.
Shannon himself used natural language to estimate the entropy of English~\cite{shannon1948}.
Recently, Deletang et al.~\cite{deletang2024compression} (ICLR 2024) argued that ``language modeling is compression,'' showing that LLMs can be viewed as general-purpose compressors that approximate the Kolmogorov complexity of text.

More recently, Ton et al.~\cite{ton2025cot} (ICML 2025) formalized Chain-of-Thought reasoning via information theory, quantifying ``information gain'' at each reasoning step---the single-agent analogue of our per-step $D_k$/$A_k$ analysis.
Nagle et al.~\cite{nagle2024ratedistortion} (NeurIPS 2024) established rate-distortion bounds for prompt compression with frozen LLMs, providing theoretical limits on how much useful information an LLM can extract from compressed input---directly relevant to our saturation analysis.
Cheng et al.~\cite{cheng2025mi} model Chain-of-Thought as a Markov chain $X \to Z \to Y$ and apply the DPI to argue that intermediate rationales preserve task-relevant information, the closest single-agent analogue of our multi-agent DPI analysis.

However, all this work focuses on \emph{single} LLM instances or single-agent reasoning chains.
The information-theoretic behavior of \emph{networks} of LLM agents---where multiple agents with different roles, capacities, and access to external tools form an information processing graph---remains unexplored.
Our work extends this perspective from single channels to \textbf{heterogeneous networks of channels}.

% --- 2.3: 組織の情報処理理論 ---
\subsection{Information Processing Theory of Organizations}

Galbraith~\cite{galbraith1974designing} proposed that organizational effectiveness depends on the fit between information processing requirements (task uncertainty) and information processing capacity (organizational design).
This view treats organizations as information processing systems---precisely the metaphor that LLM multi-agent systems literalize.

However, Galbraith's theory remained qualitative.
``Information processing capacity'' was never formally defined in terms of Shannon information.
Daft and Lengel~\cite{daft1986organizational} extended this with media richness theory, but again without mathematical formalization.

Malone~\cite{malone1987coordination} provided a more formal treatment, modeling coordination structures (product hierarchies, functional hierarchies, markets) using queuing theory to analyze tradeoffs between production costs, coordination costs, and vulnerability costs.
Computational Organization Theory~\cite{carley2002computational} extended this with agent-based simulations (e.g., ORGAHEAD, CONSTRUCT-O), demonstrating computationally that network topology determines information diffusion and task performance---but using rule-based agents, not LLMs.

LLM agents provide the first opportunity to operationalize Galbraith's theory quantitatively:
``information processing capacity'' maps to channel capacity,
``information processing requirements'' maps to input entropy,
and ``fit'' maps to the rate-distortion trade-off.
Our framework makes these mappings precise.


% ============================================================
\section{Theoretical Framework}
\label{sec:theory}
% ============================================================

% --- 3.1: 情報処理グラフモデル ---
\subsection{Multi-Agent Organizations as Information Processing Graphs}

\begin{definition}[Information Processing Graph]
A multi-agent organization is a directed acyclic graph $\mathcal{G} = (V, E, f, C)$ where:
\begin{itemize}
    \item $V = \{v_1, \ldots, v_n\}$ is the set of agent nodes;
    \item $E \subseteq V \times V$ is the set of directed communication edges;
    \item $f_i: \mathcal{X}_i \to \mathcal{Y}_i$ is the information processing function of agent $v_i$ (the LLM generation process);
    \item $C_i \in \mathbb{R}^+$ is the effective channel capacity of agent $v_i$.
\end{itemize}
\end{definition}

Each agent $v_i$ receives input $X_i$ (the concatenation of upstream agents' outputs, system prompt, and task context) and produces output $Y_i$ (the generated text).
The stochastic mapping $P(Y_i | X_i)$ is determined by the LLM's conditional generation process at temperature $\tau$.

The effective channel capacity $C_i$ is bounded above by the model's context window size, but is typically much smaller in practice due to attention degradation over long contexts~\cite{liu2024lost}.
We operationalize $C_i$ empirically through the saturation analysis described in Section~\ref{sec:metrics}.

% --- 3.2: LLMチャネルの3モード ---
\subsection{Three Operating Modes of LLM Channels}
\label{subsec:modes}

Classical communication channels are characterized by a fixed noise model.
LLM agents, however, exhibit \textbf{three qualitatively distinct operating modes}, determined by the agent's role and the ratio of input volume to effective capacity:

\begin{definition}[Channel Operating Modes]
For an LLM agent with input $X_i$ and output $Y_i$:
\begin{enumerate}
    \item \textbf{Preservation mode} ($|Y_i| \approx |X_i|$): The agent relays information with minimal transformation. Example: a coordinator forwarding instructions.
    \item \textbf{Compression mode} ($|Y_i| \ll |X_i|$): The agent filters and summarizes, discarding information deemed irrelevant. Example: a reviewer producing a short assessment from a long proposal.
    \item \textbf{Augmentation mode} ($|Y_i| > |X_i|$): The agent generates \emph{more} information than received, drawing on learned knowledge or external tools (e.g., web search). Example: a CEO synthesizing a comprehensive proposal from brief reports.
\end{enumerate}
\end{definition}

Our preliminary analysis of OrgBench data (Section~\ref{sec:preliminary}) reveals clear mode assignments:
the CEO consistently operates in augmentation mode (CR $= 1.85$, $D_k$ small, $A_k$ large);
the TL operates in compression mode (CR $= 0.35$, $D_k \approx 0.65$, $A_k \approx 0$);
the Researcher varies between augmentation (CR $= 1.76$ in deep configurations with web search) and saturation (CR $= 0.07$ in homo\_haiku where input exceeds 15{,}000 tokens);
and the CFO/CDO in mesh configurations exhibit saturation (CR $= 0.06$, $D_k \gg 0$).
These modes are strongly predicted by (a)~the agent's organizational role and (b)~the input-to-capacity ratio $|X_i| / C_i$.

% --- 3.3: DPIとその限界 ---
\subsection{The Data Processing Inequality and Its Limits}
\label{subsec:dpi}

\begin{proposition}[Classical DPI Applied to Agent Chains]
For a Markov chain $X \to Y_1 \to Y_2 \to \cdots \to Y_n$ where each $Y_k$ depends only on $Y_{k-1}$:
\begin{equation}
    I(X; Y_n) \leq I(X; Y_{n-1}) \leq \cdots \leq I(X; Y_1) \leq H(X)
    \label{eq:dpi}
\end{equation}
where $I(\cdot; \cdot)$ denotes mutual information and $H(\cdot)$ denotes entropy.
\end{proposition}

In a deep hierarchy $\text{CEO} \to \text{TL} \to \text{Researcher} \to \text{Writer} \to \text{CFO/CDO} \to \text{CEO}$, the DPI predicts that each relay step can only \emph{lose} information about the original task specification $X$.
This provides a theoretical basis for the ``depth penalty'' observed in OrgBench.

However, the DPI assumes a strict Markov chain where each agent's output depends \emph{only} on its input.
LLM agents violate this assumption in two ways:
\begin{enumerate}
    \item \textbf{Learned prior knowledge}: An agent in augmentation mode injects information from its training data, creating $I(X; Y_k) > I(X; Y_{k-1})$ if the injected knowledge is relevant to $X$.
    \item \textbf{External tool access}: A Researcher agent with web search access introduces entirely new information not present in the Markov chain.
\end{enumerate}

In our framework (Section~\ref{subsec:dpi}), these violations correspond to $A_k > 0$:
the classical DPI holds when $A_k = 0$ (preservation and compression modes), but is violated when $A_k > 0$ (augmentation mode).
The overall trend of the chain depends on whether distortion or augmentation dominates:
if $\sum D_k > \sum A_k$, information about $X$ decays with depth (as observed in OrgBench);
if $\sum A_k > \sum D_k$, the chain \emph{enriches} the task-relevant information (possible when strong augmentation agents like Researchers are present).

We therefore propose an \textbf{information accounting framework} for LLM agent chains.
We note upfront that this framework is not a theorem but a \emph{structured accounting identity}: its value lies not in the identity itself, but in the \emph{empirical regularities} of its terms---which we show are predictable from observable agent properties (Section~\ref{subsec:predictions}).
This approach follows recent work applying information-theoretic decompositions to LLM reasoning chains~\cite{ton2025cot, nagle2024ratedistortion}.
To make the framework tractable, we first introduce the notion of task information units.

\begin{definition}[Task Information Units]
\label{def:tiu}
Given a task input $X$, we define its \textbf{task information unit (TIU) set} $\mathcal{T}(X) = \{t_1, \ldots, t_m\}$ as the set of discrete, individually verifiable information elements in $X$.
TIUs are classified into four categories:
\emph{Requirements} (output specifications, e.g., ``include 3-year financial projections''),
\emph{Facts} (concrete claims, e.g., ``target: SMEs with 50--300 employees''),
\emph{Constraints} (boundary conditions, e.g., ``initial investment under \textyen 50M''), and
\emph{Queries} (implicit information needs, e.g., ``what is the market size?'').
\end{definition}

At each step $k$, we define the \textbf{preserved set}:
\begin{equation}
    \text{Pres}(k) = \{ t \in \mathcal{T}(X) : t \text{ is present in } Y_k \}
\end{equation}

``Present'' is assessed at three levels of fidelity:
literal (exact string match), semantic (paraphrase-equivalent), and functional (requirement addressed in spirit).
This yields the \textbf{information preservation rate}:
\begin{equation}
    P(k) = \frac{|\text{Pres}(k)|}{|\mathcal{T}(X)|}
    \label{eq:preservation}
\end{equation}

We can now give precise definitions to $D_k$ and $A_k$:

\begin{definition}[Information Distortion and Augmentation]
\label{def:dk_ak}
At step $k$, the \textbf{information distortion} $D_k$ and \textbf{information augmentation} $A_k$ are:
\begin{align}
    D_k &= \frac{|\text{Pres}(k{-}1) \setminus \text{Pres}(k)|}{|\mathcal{T}(X)|} \geq 0 \label{eq:dk} \\
    A_k &= \frac{|\text{Aug}(k)|}{|\mathcal{T}(X)|} \geq 0 \label{eq:ak}
\end{align}
where $\text{Aug}(k) = \{ u \in \text{Content}(Y_k) : u \notin \text{Content}(Y_{k-1}),\; u \text{ is relevant to } X \}$ is the set of new task-relevant information units introduced at step $k$.
\end{definition}

$D_k$ captures TIUs that were present in the previous step's output but are absent from the current step's output---information that was \emph{lost in transit}.
$A_k$ captures new task-relevant content that the agent \emph{adds}---from prior knowledge, external tools, or reasoning.
By construction, the preservation rate evolves as:
\begin{equation}
    P(k) = P(k{-}1) - D_k + A_k
    \label{eq:modified_dpi}
\end{equation}

Eq.~\ref{eq:modified_dpi} is an \textbf{accounting identity}---it holds by construction.
Its scientific content lies not in the identity itself, but in the empirical claim that $D_k$ and $A_k$ exhibit \emph{predictable regularities}: $D_k$ is approximately determined by the agent's compression ratio, and $A_k$ by the availability of external tools and the model's capacity (Section~\ref{subsec:predictions}).
The classical DPI corresponds to $A_k = 0$ for all $k$, yielding monotonic decay.

\subsubsection{Decomposition of $D_k$}

Not all information loss is equal.
We decompose $D_k$ into two qualitatively distinct components:
\begin{equation}
    D_k = D_k^{\text{filter}} + D_k^{\text{noise}}
    \label{eq:dk_decompose}
\end{equation}

\begin{itemize}
    \item $D_k^{\text{filter}}$: \textbf{Intentional filtering}---the agent selectively discards TIUs deemed irrelevant to its subtask.
    This is a feature, not a bug: a Team Lead extracting research instructions from a CEO's broad directive is performing useful compression.
    $D_k^{\text{filter}}$ need not reduce output quality.
    \item $D_k^{\text{noise}}$: \textbf{Unintentional loss}---TIUs that the agent \emph{should} have preserved but failed to, due to attention limitations, context truncation, or processing errors.
    $D_k^{\text{noise}}$ always reduces output quality.
\end{itemize}

The distinction is operationally testable.
If an agent in compression mode selectively retains task-critical TIUs (e.g., numerical targets, key constraints) while discarding background context, $D_k^{\text{filter}}$ dominates.
If TIU survival is uncorrelated with task relevance---particularly under saturation---$D_k^{\text{noise}}$ dominates.
Our tracer experiment (Section~\ref{sec:experiment}) is designed to distinguish these cases.

Our preliminary finding that TL compression ratio correlates positively with quality ($r = +0.169$, $p = 0.016$) suggests that TL filtering is \emph{suboptimally aggressive}: $D_k^{\text{noise}}$ is non-negligible, and higher retention (higher CR) yields better quality.

\subsubsection{Sources of $A_k$}

Augmentation arises from three sources:
\begin{equation}
    A_k = A_k^{\text{prior}} + A_k^{\text{tool}} + A_k^{\text{reason}}
\end{equation}

\begin{itemize}
    \item $A_k^{\text{prior}}$: Knowledge from the LLM's training data (e.g., the CEO enriching a brief market overview with industry benchmarks).
    \item $A_k^{\text{tool}}$: Information from external tools (e.g., the Researcher injecting web search results).
    \item $A_k^{\text{reason}}$: Novel inferences drawn from input (e.g., the Writer identifying a risk not mentioned by upstream agents).
\end{itemize}

$A_k^{\text{tool}}$ is directly measurable from tool call logs.
$A_k^{\text{prior}}$ can be estimated as new entities in the output that are absent from both the input and tool results.
$A_k^{\text{reason}}$ is difficult to isolate automatically and is treated as residual in this work.

Our finding that CEO input volume positively predicts quality ($r = +0.207$, $p < 0.001$) suggests that $A_k^{\text{prior}}$ is an increasing function of input volume: a richer input stimulates the CEO to draw on more relevant prior knowledge.

\subsubsection{Mode--$D_k$/$A_k$ Correspondence}

The three operating modes (Definition~2) map cleanly onto the $D_k$/$A_k$ framework:

\begin{table}[t]
\centering
\caption{Operating modes and their $D_k$/$A_k$ signatures.}
\label{tab:mode_dk_ak}
\small
\begin{tabular}{lcccc}
\toprule
\textbf{Mode} & $D_k$ & $A_k$ & \textbf{Net} & \textbf{Example} \\
\midrule
Preservation & $\approx 0$ & $\approx 0$ & $\approx 0$ & Coordinator relay \\
Compression  & $> 0$       & $\approx 0$ & $< 0$       & TL (CR$=$0.35) \\
Augmentation & small       & large       & $> 0$       & CEO (CR$=$1.85) \\
Saturation   & $\gg 0$     & $\approx 0$ & $\ll 0$     & CFO/mesh (CR$=$0.06) \\
\bottomrule
\end{tabular}
\end{table}

Saturation (Table~\ref{tab:mode_dk_ak}, bottom row) is an extreme form of compression where $D_k^{\text{noise}} \gg D_k^{\text{filter}}$: the agent cannot process its input and drops information non-selectively.

% --- 3.4: チャネル飽和 ---
\subsection{Channel Saturation}
\label{subsec:saturation}

\begin{definition}[Channel Saturation]
An LLM agent $v_i$ is \textbf{saturated} when its input volume $|X_i|$ exceeds its effective capacity $C_i$, causing a discontinuous drop in output quality:
\begin{equation}
    \text{CR}(v_i) = \frac{|Y_i|}{|X_i|} \to \epsilon \quad \text{as} \quad |X_i| \gg C_i
\end{equation}
where $\epsilon \ll 1$.
In the saturated regime, the agent's output contains a vanishingly small fraction of the input information.
\end{definition}

% [TODO: 飽和閾値の理論的・実証的同定]
% 予備分析からの示唆:
% - CFO (GPT-4o-mini): 通常入力 ~1,600 tokens → CR=0.35
%   mesh入力 ~7,500 tokens → CR=0.06 → 飽和閾値は 3,000-5,000 tokens?
% - CEO (Haiku): 入力 3,000-12,000 tokens → CR=1.8-1.9 → 飽和せず
%   → Haikuの飽和閾値は 12,000+ tokens?
% - 飽和閾値 ≈ f(model_capability, task_complexity)?

Our preliminary data suggests that channel saturation occurs at a model-dependent threshold.
For GPT-4o-mini agents, saturation appears between 3,000--5,000 input tokens; for Claude Haiku 4.5, the threshold exceeds 12,000 tokens.
This has direct design implications: \textbf{mesh topologies that route all information to all agents will saturate weak agents}, negating the theoretical advantage of path redundancy.

% --- 3.5: ボトルネック原理 ---
\subsection{The Bottleneck Principle}
\label{subsec:bottleneck}

Drawing on the max-flow min-cut theorem from network flow theory, one might expect that overall quality is limited by the minimum-capacity agent.
However, our empirical analysis of 250 runs reveals a more nuanced picture.

\begin{proposition}[Depth-Distortion Principle]
The information throughput of a multi-agent organization $\mathcal{G}$ is primarily determined by the \textbf{number of relay steps on the critical path}---the sequence of agents through which information must pass from task input to final output.
Each relay step introduces a roughly constant rate of information loss:
\begin{equation}
    I(X; Y_n) \approx I(X; Y_0) \cdot \alpha^d
    \label{eq:depth_distortion}
\end{equation}
where $d$ is the critical path depth and $\alpha < 1$ is the per-step retention rate.
\end{proposition}

We validate this through correlation analysis (Section~\ref{sec:preliminary}).
Critical path depth is the strongest predictor of output quality ($r = -0.255$, $p < 0.0001$), outperforming both min-CR ($r = +0.093$, $p = 0.14$, n.s.) and mean-CR ($r = +0.039$, $p = 0.54$, n.s.).

Crucially, agents \emph{not on the critical path} (e.g., reviewers) do not contribute to depth-distortion even when saturated.
In flat\_mesh, the CFO and CDO exhibit CR $= 0.06$ (severe saturation), yet this configuration achieves the \emph{highest} quality (OQ $= 3.30$).
This is because reviewers provide feedback but do not relay the main information flow.

This principle predicts that:
\begin{itemize}
    \item Adding an intermediate relay layer (TL with CR $= 0.35$) reduces quality by introducing one additional distortion step, regardless of other agents' capacities.
    \item The flat\_hub configuration outperforms deep hierarchies because it \emph{eliminates the TL relay step}, reducing critical path depth from 4 to 3.
    \item Mesh topologies can improve quality by increasing the CEO's input volume ($r = +0.207$, $p < 0.001$), but only if the \emph{critical-path} agents have sufficient capacity.
    \item Reviewer saturation is tolerable; critical-path saturation is not.
\end{itemize}

% --- 3.6: トポロジー別の理論的予測 ---
\subsection{Topology-Specific Predictions}
\label{subsec:predictions}

We derive predictions for the topologies present in OrgBench:

\begin{table}[t]
\centering
\caption{Theoretical predictions by topology.}
\label{tab:predictions}
\small
\begin{tabular}{p{1.5cm}p{2.2cm}p{1.5cm}p{2cm}}
\toprule
\textbf{Topology} & \textbf{Info. Path Length} & \textbf{Bottleneck} & \textbf{Prediction} \\
\midrule
Chain (deep) & $n$ (max) & TL (CR=0.35) & High distortion \\
Star (flat\_hub) & 2 (max) & CEO only & Low distortion \\
Mesh (flat\_mesh) & 1 (direct) & CFO/CDO if saturated & Low distortion \emph{if} capacity sufficient \\
\bottomrule
\end{tabular}
\end{table}

% --- 3.7: 検証可能な予測（反直感的予測を含む） ---
\subsection{Testable Predictions}
\label{subsec:testable}

Our framework generates six testable predictions, including three that are counterintuitive and would not be derived without the information-flow analysis:

\textbf{H1 (Depth-Distortion, confirmatory):}
Tracer survival rate decays with critical path depth $d$.
Expected: $\text{TSR}(d) \propto \alpha^d$ with $\alpha < 1$.

\textbf{H2 (TL as cause, confirmatory):}
Replacing the TL with a pass-through relay (system prompt: ``forward all information without filtering'') eliminates or significantly reduces the depth penalty.
If depth=4 with pass-through TL $\approx$ depth=3 without TL, then the TL's filtering---not depth per se---is the primary cause of quality loss.
If depth=4 with pass-through TL $<$ depth=3, then information distortion occurs even in pure relay, implicating the LLM channel itself.

\textbf{H3 (Capacity-dependent depth penalty, \emph{counterintuitive}):}
The depth penalty is \emph{model-dependent}: high-capacity models exhibit smaller $D_k$ per step, reducing the decay rate $\alpha$.
\emph{Prediction}: For sufficiently capable models, deep hierarchies can match or exceed shallow ones, because deeper structures decompose the task into smaller subtasks that each agent can process more effectively.
\emph{Evidence}: homo\_haiku achieves depth=4, OQ$=$3.22---matching flat\_hub (depth=3, OQ$=$3.22)---because Haiku's TL retains CR$=$0.87 vs.\ GPT-4o-mini's CR$=$0.35.
In contrast, homo\_gpt (depth=4, all GPT-4o-mini) achieves only OQ$=$2.56, confirming that the depth penalty is severe for low-capacity models but vanishes for high-capacity ones.

\textbf{H4 (Mesh harms weak agents, \emph{counterintuitive}):}
Mesh topology \emph{reduces} quality when the organization contains low-capacity agents, because mesh routing saturates them.
\emph{Prediction}: In heterogeneous configurations (strong + weak models), hub-and-spoke outperforms mesh; in homogeneous high-capacity configurations, mesh outperforms hub.

\textbf{H5 (Review is dimension-selective):}
Review gates function as error-correction channels, primarily reducing $D_k^{\text{noise}}$ in dimensions where specialist reviewers have domain expertise.
\emph{Prediction}: Review protocols improve quality dimensions aligned with reviewer expertise (Technical Depth for CDO review, Feasibility for integrated review) more than dimensions orthogonal to it.
\emph{Preliminary evidence}: Comparing anchor (review) vs.\ deep\_noreview, the largest gains are in Technical Depth ($\Delta = +0.34$) and Feasibility ($\Delta = +0.20$), while Novelty ($\Delta = +0.08$) and Financial Rigor ($\Delta = -0.02$) show negligible change.

\textbf{H6 (Augmentation is input-dependent):}
$A_k^{\text{prior}}$ for augmentation-mode agents (CEO) is an increasing function of input volume.
\emph{Prediction}: CEO quality improves with more input information, with no saturation within current context window limits.


% ============================================================
\section{Information-Theoretic Metrics}
\label{sec:metrics}
% ============================================================

We define six metrics to quantify information flow through agent organizations.

\subsection{Token-Level Metrics (Model-Free)}

\textbf{M1: Compression Ratio (CR).}
\begin{equation}
    \text{CR}(v_i) = \frac{|Y_i|_{\text{tokens}}}{|X_i|_{\text{tokens}}}
\end{equation}
Measures the token-level expansion or contraction at each agent.
$\text{CR} > 1$ indicates augmentation; $\text{CR} \ll 1$ with high input indicates saturation.

% [STATUS: ✅ 既存データで計算済み（予備分析）]

\textbf{M2: Information Density (ID).}
\begin{equation}
    \text{ID}(Y_i) = \frac{|\text{unique n-grams}(Y_i)|}{|Y_i|_{\text{tokens}}}
\end{equation}
Measures lexical diversity. Low ID indicates repetitive or formulaic output, a symptom of saturation.

% [TODO: 実装 — n=3 (trigram) を標準とする。既存output.mdから計算可能]

\subsection{Semantic Metrics (Embedding-Based)}

\textbf{M3: Semantic Similarity Decay (SSD).}
\begin{equation}
    \text{SSD}(k) = \cos\bigl(\text{emb}(X_0),\, \text{emb}(Y_k)\bigr)
\end{equation}
where $X_0$ is the original task prompt and $Y_k$ is the output at step $k$.
Measures how far the output has drifted from the original task specification.

% [TODO: 実装 — text-embedding-3-small を使用。全文テキストが必要 → orchestrator修正後]

\textbf{M4: Cross-Step Mutual Information Estimate (CSMI).}
\begin{equation}
    \hat{I}(Y_{k-1}; Y_k) \approx g\bigl(\cos(\text{emb}(Y_{k-1}), \text{emb}(Y_k))\bigr)
\end{equation}
Embedding-based approximation of mutual information between consecutive steps.

% [TODO: 推定手法の選定 — KSG estimator vs MINE vs 簡易cosine-based proxy]
% 実用上はcosine similarityをMIのproxyとして使い、
% 厳密なMI推定はablationとして実施するのが現実的か？

\subsection{Content Preservation Metrics}

\textbf{M5: Entity Preservation Rate (EPR).}
\begin{equation}
    \text{EPR}(X_i, Y_i) = \frac{|\text{NER}(X_i) \cap \text{NER}(Y_i)|}{|\text{NER}(X_i)|}
\end{equation}
Fraction of named entities in the input that survive to the output.

% [TODO: 実装 — spaCy日本語モデル or multilingual NER]
% 課題: 日本語テキストのNER精度。英語に翻訳してからNERする方が精度が高い可能性

\textbf{M6: Numerical Fidelity (NF).}
\begin{equation}
    \text{NF}(X_i, Y_i) = \frac{|\text{nums}(X_i) \cap \text{nums}(Y_i)|}{|\text{nums}(X_i)|}
\end{equation}
Fraction of numerical values that are accurately preserved.

% [TODO: 実装 — regex-based数値抽出。「1兆3,412億円」等の日本語数値表現への対応が必要]

\subsection{Controlled Information Injection}

\textbf{M7: Tracer Survival Rate (TSR).}
We inject $K$ synthetic ``tracer'' facts into the task prompt (fictitious company names, fabricated statistics, specific constraints) and measure what fraction survives to the final output:
\begin{equation}
    \text{TSR} = \frac{1}{K} \sum_{j=1}^{K} \mathbb{1}[\text{tracer}_j \in Y_{\text{final}}]
\end{equation}

% [TODO: 設計と実装]
% - トレーサーの種類: 架空固有名詞、架空数値、架空事実、制約条件
% - 各テーマに10個のトレーサーを埋め込み
% - 検出方法: exact match + fuzzy match (数値の丸め等に対応)
% - 各ステップでの検出率を追跡 → 生存曲線を構築

This metric provides the most direct test of the DPI: if information is monotonically lost across relay steps, TSR should decrease with chain depth.
Unlike M3--M6, TSR is independent of embedding models and NER accuracy, providing a model-free ground truth for information preservation.


% ============================================================
\section{Preliminary Empirical Analysis}
\label{sec:preliminary}
% ============================================================

Before conducting the full information-theoretic experiment (Section~\ref{sec:experiment}), we analyze the existing OrgBench pilot data (275 runs across 11 configurations) to validate our theoretical framework and guide experimental design.

% --- 5.1: CR分析 ---
\subsection{Compression Ratio Analysis}

Table~\ref{tab:cr} presents the mean compression ratio by agent role and configuration.

\begin{table}[t]
\centering
\caption{Mean compression ratio (CR = output tokens / input tokens) by agent and configuration. Bold indicates extreme values discussed in text.}
\label{tab:cr}
\small
\begin{tabular}{lcccccc}
\toprule
\textbf{Config} & \textbf{CEO} & \textbf{TL} & \textbf{Res.} & \textbf{Wri.} & \textbf{CFO} & \textbf{CDO} \\
\midrule
anchor       & 1.85 & 0.35 & 1.61 & 0.66 & 0.35 & 0.23 \\
flat\_hub    & 1.90 & ---  & 1.10 & 0.75 & 0.33 & 0.26 \\
flat\_mesh   & 1.80 & ---  & 0.50 & 0.51 & \textbf{0.06} & \textbf{0.06} \\
deep\_noreview & 1.92 & 0.35 & 1.76 & 0.66 & ---  & ---  \\
homo\_haiku  & 1.81 & 0.87 & 0.07 & 1.17 & 0.83 & 0.83 \\
homo\_gpt    & 1.31 & 0.54 & 0.15 & 0.72 & 0.36 & 0.21 \\
\bottomrule
\end{tabular}
\end{table}

Three findings emerge:

\textbf{Finding 1: Role determines operating mode.}
The CEO consistently operates in augmentation mode (CR $\approx$ 1.85), synthesizing comprehensive proposals from partial reports.
The TL operates in compression mode (CR $\approx$ 0.35), filtering 65\% of input.
These roles are stable across configurations, suggesting that the agent's \emph{system prompt} (not the topology) determines its channel mode.

\textbf{Finding 2: Channel saturation is real and topology-dependent.}
In flat\_mesh, the CFO and CDO receive $\sim$7,500 tokens (4$\times$ their typical input) and their CR collapses to 0.06.
This is not gradual degradation but a phase transition: the agents effectively cease to produce useful output.
The saturation threshold for GPT-4o-mini appears to be 3,000--5,000 tokens for the reviewer role.

\textbf{Finding 3: Model capability ≈ channel capacity.}
In homo\_haiku (all Claude Haiku 4.5), most agents maintain CR $>$ 0.8 regardless of input volume, and TL achieves CR = 0.87 vs. 0.35 in heterogeneous configurations.
Haiku's larger effective capacity allows it to preserve more information at each step.

% --- 5.2: ボトルネック分析 ---
\subsection{Bottleneck Analysis}

\begin{table}[t]
\centering
\caption{CEO final integration step: mean input tokens vs. output quality.}
\label{tab:bottleneck}
\small
\begin{tabular}{lccc}
\toprule
\textbf{Config} & \textbf{CEO input} & \textbf{OQ} & \textbf{Bottleneck agent} \\
\midrule
flat\_mesh   & 10,591 & \textbf{3.30} & CFO/CDO (saturated) \\
homo\_haiku  & 9,459  & \textbf{3.22} & none (all high-cap.) \\
flat\_hub    & 5,331  & \textbf{3.22} & none \\
anchor       & 7,064  & 2.92          & TL (CR=0.35) \\
deep\_noreview & 4,941 & 2.68         & TL (CR=0.35) \\
homo\_gpt    & 4,759  & 2.56          & all (low-cap. model) \\
\bottomrule
\end{tabular}
\end{table}

Table~\ref{tab:bottleneck} reveals a counterintuitive result: the CEO's input volume \emph{positively} correlates with quality ($r = +0.207$, $p < 0.001$, $n = 250$).
More information reaching the CEO leads to better output, not worse.

Table~\ref{tab:correlations} presents the full correlation analysis.
The strongest predictor of output quality is \textbf{critical path depth} ($r = -0.255$, $p < 0.0001$)---configurations with fewer relay steps produce higher quality.
The second strongest is CEO input volume ($r = +0.207$, $p < 0.001$).
The TL's compression ratio also shows a significant positive correlation ($r = +0.169$, $p = 0.016$): when the TL preserves more information, quality improves.

Notably, min-CR across all agents shows a \emph{negative} correlation with quality ($r = -0.135$, $p = 0.032$).
This seemingly paradoxical result arises because the highest-quality configuration (flat\_mesh) has severely saturated reviewers (CFO/CDO at CR $= 0.06$), which pull down the min-CR but do not affect the main information flow.
When restricted to critical-path agents, the correlation becomes non-significant ($r = +0.093$, $p = 0.14$), confirming that \textbf{reviewer saturation is harmless while critical-path depth is the primary quality determinant}.

\begin{table}[t]
\centering
\caption{Correlation of candidate predictors with overall quality ($n = 250$ multi-agent runs).}
\label{tab:correlations}
\small
\begin{tabular}{lccl}
\toprule
\textbf{Predictor} & \textbf{Pearson} $r$ & $p$ & \textbf{Sig.} \\
\midrule
Critical path depth       & $-0.255$ & $< .0001$ & *** \\
CEO final input tokens    & $+0.207$ & $.0009$   & ** \\
TL compression ratio      & $+0.169$ & $.016$    & * \\
Min CR (all agents)       & $-0.135$ & $.032$    & * \\
Writer CR                 & $+0.114$ & $.071$    & n.s. \\
Min CR (critical path)    & $+0.093$ & $.141$    & n.s. \\
Mean CR (critical path)   & $+0.039$ & $.535$    & n.s. \\
\bottomrule
\end{tabular}
\end{table}

These findings support the Depth-Distortion Principle (Proposition~3) rather than a simple min-capacity bottleneck model.
The causal mechanism for the ``depth penalty'' appears to be:
(1)~the TL acts as a fixed-rate filter (CR $\approx 0.35$), discarding 65\% of input;
(2)~this reduces the information available to downstream agents;
(3)~consequently, the CEO's final integration step receives less raw material, producing a lower-quality synthesis.
Eliminating the TL (flat configurations) removes this filter, allowing more information to reach the CEO and yielding higher quality.


% ============================================================
\section{Experimental Design}
\label{sec:experiment}
% ============================================================

% [TODO: 全セクション — Phase 1追加実験の設計]

Building on the preliminary analysis, we design three targeted experiments to validate our theoretical predictions.

\subsection{Experiment A: Full-Trace Re-Run}

% [TODO: orchestrator.py修正後に実行]
% 目的: 各ステップの全文テキストを取得し、M2-M6を計算可能にする
% 構成: anchor, flat_hub, flat_mesh, deep_noreview (4構成)
% テーマ: 5テーマ × 3反復 = 60回
% 推定コスト: ~$2.4

We re-run 4 representative configurations (anchor, flat\_hub, flat\_mesh, deep\_noreview) with full-text step traces, enabling computation of all six metrics across every agent step.

\subsection{Experiment B: Chain Depth Variation}

% [TODO: chain_2〜6のYAMLテンプレート作成、実行]
% 目的: Information Distortion Curveの構築
% 構成: chain長 2/3/4/5/6 (5構成)
% テーマ: 5テーマ × 3反復 = 75回
% 推定コスト: ~$2.3

To directly test the DPI prediction that information fidelity decays with chain depth, we construct 5 configurations with chain lengths 2 through 6, holding all other variables constant.

\subsection{Experiment C: Tracer Injection}

% [TODO: トレーサー付きテーマYAML作成、trace.py実装、実行]
% 目的: 情報保存の直接測定
% 構成: 4構成 × 3テーマ(トレーサー付き) × 3反復 = 36回
% 推定コスト: ~$1.4

We inject 10 synthetic tracer facts per theme and track their survival across agent steps, providing a model-free ground truth for information preservation.


% ============================================================
\section{Results}
\label{sec:results}
% ============================================================

% [TODO: 実験完了後に記述]

\subsection{Information Distortion Curve}
% [TODO: Experiment B のデータから]
% - chain長 d vs 各情報指標（M2-M7）のプロット
% - 減衰モデル（線形/指数/対数）のフィッティング、AIC/BICで選択
% - 理論予測（DPIによる単調減少）との一致度

\subsection{Topology-Dependent Information Preservation}
% [TODO: Experiment A のデータから]
% - 構成別のM2-M7指標の比較
% - Kruskal-Wallis検定でトポロジー間の差異を検定
% - Table~\ref{tab:predictions}の理論予測 vs 実測

\subsection{Channel Saturation Threshold}
% [TODO: Experiment A のデータから]
% - 入力トークン数 vs CR の散布図（エージェント別）
% - breakpoint回帰で飽和閾値を同定
% - モデル別（GPT-4o-mini vs Haiku）の閾値比較

\subsection{Tracer Survival Analysis}
% [TODO: Experiment C のデータから]
% - 各ステップでのトレーサー検出率
% - 構成別×トレーサー種別のヒートマップ
% - chain長 vs TSR の相関

\subsection{Predictive Power of Information Metrics}
% [TODO: 全実験データを統合]
% - 回帰分析: OQ ~ M1 + M2 + ... + M7
% - 情報指標 vs 構造指標（次数、経路長）の予測力比較
% - どの指標が最も品質を予測するか


% ============================================================
\section{Discussion}
\label{sec:discussion}
% ============================================================

\subsection{Do LLM Agent Chains Obey the Data Processing Inequality?}

% [TODO: 実験結果に基づく議論]
% 予想される結論:
% - Pure relay/compression ステップではDPIが成立（TLでの情報損失）
% - Augmentation ステップ（CEO, Researcher）ではDPIが破れる
% - 全体としてはDPIの「修正版」（Eq. 2）が記述的に有効
% - 重要な知見: 情報の「量」だけでなく「質」が問題。
%   TLは65%をカットするが、残り35%が適切な情報ならば問題ない。
%   CFO/CDOが94%をカットするのは、処理しきれない（飽和）ため。

\subsection{Channel Saturation as the Primary Mechanism}

% [TODO: 実験結果に基づく議論]
% 「なぜ深い階層が悪いか」の最終回答:
% - 仮説(a) information distortion: 部分的に正しい（TLでの情報損失）
% - 仮説(b) over-filtering: 部分的に正しい（TLのCR=0.35）
% - 仮説(c) resource competition: 証拠不十分
% - 新仮説: ボトルネック原理。深い階層は低容量ノード（TL）を挿入するため、
%   全体の情報スループットが低下する。

\subsection{Practical Design Guidelines}

% 情報理論に基づく設計原則（Phase 0の経験則を理論的に裏付け）:

% Guideline 1: 最弱リンクを排除せよ
% - 組織全体の品質は min(C_i) で制約される
% - 低容量エージェント（小モデルの中間層）を挿入するくらいなら省略した方が良い
% → flat_hub > anchor の理論的説明

% Guideline 2: 受信側の容量を考慮してトポロジーを選べ
% - meshは理論的には最適だが、受信側が飽和すれば逆効果
% - 全エージェントの容量が十分なら mesh > hub > chain
% - 容量が不均一なら hub（情報集約）+ 高容量ハブが最適

% Guideline 3: 情報増幅ノードはボトルネックにならない
% - CEO（CR>1）は入力が増えても品質が落ちない（容量が十分な限り）
% - 設計上、augmentationモードのノードをハブに配置すべき

% [TODO: 上記をフォーマルに記述]

\subsection{Limitations}

% [TODO: 正直に限界を記述]
% 1. Compression Ratioは情報量の粗い代理指標。トークン数 ≠ 情報量
% 2. Embedding-based MIは近似。厳密なMI計算は高次元テキストでは困難
% 3. タスクがビジネス提案のみ。他ドメインへの一般化は未検証
% 4. モデル3種のみ（Haiku, GPT-4o-mini, Gemini Flash）。他モデルでの検証必要
% 5. チャネル飽和閾値はモデル×役割×タスク依存。普遍的閾値は存在しない可能性
% 6. 理論モデルは定性的予測に留まる。定量的予測（「CRが0.3なら品質は2.8」）はまだ不可能


% ============================================================
\section{Conclusion}
\label{sec:conclusion}
% ============================================================

% [TODO: 実験完了後に記述]

% 主要貢献の要約:
% 1. LLMマルチエージェント組織を情報処理グラフとしてモデル化する理論的枠組み
% 2. LLMチャネルの3モード（保存/圧縮/増幅）の発見と特徴づけ
% 3. チャネル飽和現象の同定 — 「深い階層が悪い」理由の因果的説明
% 4. ボトルネック原理 — 設計指針の理論的根拠
% 5. 6つの情報理論的指標による実証検証

% Future work:
% - Phase 2（動的適応）への接続:
%   「information processing demand > channel capacity のとき構造を変更すべき」
%   という理論的トリガーの導出
% - チャネル容量の事前推定手法（実験なしで飽和閾値を予測する方法）


% ============================================================
% References
% ============================================================

\bibliographystyle{IEEEtran}

\begin{thebibliography}{20}

\bibitem{wu2023autogen}
Q.~Wu et al.,
``AutoGen: Enabling next-gen LLM applications via multi-agent conversation,''
in \emph{Proc.\ COLM}, 2024.

\bibitem{hong2023metagpt}
S.~Hong et al.,
``MetaGPT: Meta programming for a multi-agent collaborative framework,''
in \emph{Proc.\ ICLR}, 2024.

\bibitem{qian2024chatdev}
C.~Qian et al.,
``ChatDev: Communicative agents for software development,''
in \emph{Proc.\ ACL}, 2024.

\bibitem{qian2025macnet}
C.~Qian et al.,
``Scaling large language model-based multi-agent collaboration,''
in \emph{Proc.\ ICLR}, 2025.

\bibitem{yao2025gdesigner}
Y.~Yao et al.,
``G-Designer: Architecting multi-agent communication topologies via graph neural networks,''
in \emph{Proc.\ ICML}, 2025.

\bibitem{multiagentbench2025}
% [TODO: 正確な著者情報]
``MultiAgentBench: Evaluating the collaboration and competition of LLM agents,''
in \emph{Proc.\ ACL}, 2025.

\bibitem{shinoda2026orgbench}
K.~Shinoda,
``OrgBench: How information flow design patterns affect collective decision quality in LLM multi-agent systems,''
2026. [Manuscript in preparation]

\bibitem{shannon1948}
C.~E. Shannon,
``A mathematical theory of communication,''
\emph{The Bell System Technical Journal}, vol.~27, pp.~379--423, 623--656, 1948.

\bibitem{deletang2024compression}
G.~Deletang et al.,
``Language modeling is compression,''
in \emph{Proc.\ ICLR}, 2024.

\bibitem{galbraith1974designing}
J.~R. Galbraith,
``Organization design: An information processing view,''
\emph{Interfaces}, vol.~4, no.~3, pp.~28--36, 1974.

\bibitem{daft1986organizational}
R.~L. Daft and R.~H. Lengel,
``Organizational information requirements, media richness and structural design,''
\emph{Management Science}, vol.~32, no.~5, pp.~554--571, 1986.

\bibitem{liu2024lost}
N.~F. Liu et al.,
``Lost in the middle: How language models use long contexts,''
\emph{TACL}, 2024.

\bibitem{ton2025cot}
J.-H.~Ton, M.~Taufiq, and Y.~Liu,
``Understanding chain-of-thought in LLMs through information theory,''
in \emph{Proc.\ ICML}, 2025.

\bibitem{nagle2024ratedistortion}
J.~Nagle et al.,
``Fundamental limits of prompt compression: A rate-distortion framework for black-box language models,''
in \emph{Proc.\ NeurIPS}, 2024.

\bibitem{cheng2025mi}
Z.~Cheng et al.,
``Quantifying information gain and redundancy in multi-turn LLM interactions,''
\emph{arXiv preprint}, 2025.

\bibitem{agentprune2025}
H.~Zhang et al.,
``AgentPrune: An economical communication pipeline for LLM-based multi-agent systems,''
in \emph{Proc.\ ICLR}, 2025.

\bibitem{agentdropout2025}
S.~Wang et al.,
``AgentDropout: Dynamic agent elimination for token-efficient and high-performance multi-agent communication,''
in \emph{Proc.\ ACL}, 2025.

\bibitem{shen2025infoprop}
Y.~Shen et al.,
``Understanding the information propagation effects of communication topologies in LLM-based multi-agent systems,''
\emph{arXiv:2505.23352}, 2025.

\bibitem{cemri2025fail}
M.~Cemri et al.,
``Why do multi-agent LLM systems fail?''
in \emph{Proc.\ NeurIPS}, 2025.

\bibitem{malone1987coordination}
T.~W. Malone,
``Modeling coordination in organizations and markets,''
\emph{Management Science}, vol.~33, no.~10, pp.~1317--1332, 1987.

\bibitem{carley2002computational}
K.~M. Carley,
``Computational organization science: A new frontier,''
\emph{PNAS}, vol.~99, suppl.~3, pp.~7257--7262, 2002.

\end{thebibliography}

\end{document}
